\begin{lemma}[Selection of bad scales]
Fix an integer $d\ge1$, put $n=d+1$, write $G=B_{3/4}$ for the open ball of
radius $3/4$ about the origin of the space-time $\ell^2(\R^d\times\R)$ and
$D(z)=\operatorname{dist}(z,\partial G)$, and let $b_d$ be the family flatness
modulus of \noderef{FlatnessModulus}. Suppose that the uniform power bound on
$b_d$ fails for every positive exponent: for every $B>0$, every
$\gamma\in(0,1]$ and every $s_0\in(0,1/8]$ there is $s\in(0,s_0]$ with
\[
  b_d(s)>Bs^\gamma .
\]
Then there are real sequences $\gamma_i,s_i,r_i,\lambda_i$, functions $U_i$ on
the space-time and centers $z_i$ such that
$(\gamma_i,s_i,r_i,\lambda_i,U_i,z_i)_{i\ge0}$ is selected bad-scale data in
dimension $d$ in the sense of \noderef{IsBadScaleData}.
\end{lemma}
\begin{proof}
Throughout, $\mathcal U$ is the normalized family of
\noderef{NormalizedInfFamily}, $\beta_U(z,r)$ is as in \noderef{BetaFlatness},
and by \noderef{FlatnessModulus}
\begin{equation}\label{eq:bs-bdef}
  b_d(s)=\sup\bigl\{\beta_U(z,r):U\in\mathcal U,\ z\in G,\ 0<r\le sD(z)\bigr\}.
\end{equation}
By \noderef{BetaFlatness} each $\beta_U(z,r)$ is the supremum of a subset of
$(0,1)$, so $0\le\beta_U(z,r)\le1$; hence the set in \eqref{eq:bs-bdef} is
bounded above by $1$ and $0\le b_d(s)\le1$. What this establishes is boundedness
above, and nothing more: we do \emph{not} claim here that the set in
\eqref{eq:bs-bdef} is nonempty. Nonemptiness is needed only where the
\emph{approximation property} of a supremum is invoked, which happens exactly
three times: for the scalar supremum $M$ in Step 1, and in Step 2 for
\eqref{eq:bs-bdef} itself and for $\beta_{U_i}(z_i,r_i)$. It is established at
each of those three points of use, in Step 2 from the strict positivity
\eqref{eq:bs-bpos}. (The scalar supremum $T$ of Step 1 is used only as an upper
bound, together with $T\ge0$, and its defining set is visibly nonempty because
$1/N\le1/8$.) We use three consequences repeatedly: $\beta_U(z,r)\le
b_d(s)$ whenever $(U,z,r)$ is admissible for $s$ as in \eqref{eq:bs-bdef}, which
needs only the boundedness just proved; every $\lambda$ with
$\beta_U(z,r)<\lambda<1$ is good for $U$ at $(z,r)$, i.e.\ admits a cylinder fit
in the sense of \noderef{IsCylinderFit}, since otherwise $\lambda$ would lie in
the set whose supremum is $\beta_U(z,r)$; and $D$ is $1$-Lipschitz, being a
distance to a fixed set.

We also record that for $z\in G$
\begin{equation}\label{eq:bs-D}
  D(z)=\tfrac34-|z|\in(0,\tfrac34],
\end{equation}
because $\partial G$ is the sphere of radius $3/4$, which is nonempty since the
space-time has dimension $d+1\ge2$. Indeed the point $\tfrac34|z|^{-1}z$, or any
point of the sphere when $z=0$, lies on the sphere at distance exactly
$\tfrac34-|z|$ from $z$; and no point $w$ of the sphere is nearer, because
$|z-w|\ge|w|-|z|=\tfrac34-|z|$.

\emph{Step 1: the running maximum.}
Put $\gamma_i=1/(i+1)$ for $i\ge0$, so $0<\gamma_i\le1$ and $\gamma_i\to0$. Fix
$i$ and let
\[
  g(s)=\frac{b_d(s)}{s^{\gamma_i}}\qquad(0<s\le1/8),
\]
a nonnegative function. Let $N$ be an integer with $N\ge\max(8,i+1)$, so that
$1/N\le1/8$.

The quantity $T=\sup\{g(s):1/N\le s\le1/8\}$ is finite: on that interval
$b_d(s)\le1$ and $s^{-\gamma_i}\le N^{\gamma_i}$, so $g\le N^{\gamma_i}$. Also
$T\ge0$. Apply the hypothesis with $B=2T+2>0$, the exponent $\gamma_i$ and
$s_0=1/N$: there is $\sigma\in(0,1/N]$ with $b_d(\sigma)>(2T+2)\sigma^{\gamma_i}$,
that is
\[
  g(\sigma)>2T+2 .
\]
In particular $g(\sigma)>T$, so $\sigma\notin[1/N,1/8]$ and therefore
\begin{equation}\label{eq:bs-sigma}
  0<\sigma<\tfrac1N .
\end{equation}

Let $M=\sup\{g(s):\sigma\le s\le1/8\}$, again finite because $b_d\le1$ and
$s^{-\gamma_i}\le\sigma^{-\gamma_i}$ there. Its defining set is nonempty, and
the witness is $\sigma$ itself: by \eqref{eq:bs-sigma} and $N\ge8$ we have
$\sigma<1/N\le1/8$, so $\sigma\in[\sigma,1/8]$ and $g(\sigma)$ belongs to the
set. Hence $M\ge g(\sigma)>2$, and $M$ is the supremum of a nonempty set bounded
above. Since $M/2<M$, the approximation property of the supremum is available
and gives $s_i\in[\sigma,1/8]$ with
\begin{equation}\label{eq:bs-approxmax}
  g(s_i)>\tfrac12M .
\end{equation}
Two consequences. First, every $s\in[s_i,1/8]$ lies in $[\sigma,1/8]$, so
$g(s)\le M<2g(s_i)$, i.e.
\begin{equation}\label{eq:bs-running}
  \frac{b_d(s)}{s^{\gamma_i}}\le2\frac{b_d(s_i)}{s_i^{\gamma_i}}
  \qquad(s_i\le s\le1/8).
\end{equation}
Second, $g(s_i)>\tfrac12M\ge\tfrac12g(\sigma)>T$, while every $s\in[1/N,1/8]$
has $g(s)\le T$; so $s_i\notin[1/N,1/8]$, and as $s_i\in[\sigma,1/8]$ with
$\sigma<1/N$ by \eqref{eq:bs-sigma}, we conclude
\begin{equation}\label{eq:bs-si}
  0<s_i<\tfrac1N\le\frac1{i+1},\qquad s_i\le\tfrac18 .
\end{equation}
Hence $s_i\to0$. Also $g(s_i)>\tfrac12M>1>0$, so
\begin{equation}\label{eq:bs-bpos}
  b_d(s_i)>0 .
\end{equation}

\emph{Step 2: the bad triple and the bad parameter.}
Fix $i$. The approximation property of a supremum is available only for a
nonempty set, and here strict positivity supplies that. By
\noderef{FlatnessModulus} the value $b_d(s_i)$ is $0$ when the set in
\eqref{eq:bs-bdef} is empty, whereas $b_d(s_i)>0$ by \eqref{eq:bs-bpos}; hence
that set is nonempty, and being bounded above by $1$ as recorded above it has a
genuine supremum. Since $b_d(s_i)/2<b_d(s_i)$, again by \eqref{eq:bs-bpos}, the
approximation property at $s=s_i$ gives $U_i\in\mathcal U$, $z_i\in G$ and $r_i$
with
\begin{equation}\label{eq:bs-triple}
  0<r_i\le s_iD(z_i),\qquad \beta_{U_i}(z_i,r_i)>\tfrac12b_d(s_i)>0 .
\end{equation}
The same argument runs one level down. By \noderef{BetaFlatness} the quantity
$\beta_{U_i}(z_i,r_i)$ is the supremum of the set of parameters in $(0,1)$ that
are not good for $U_i$ at $(z_i,r_i)$, \emph{with the value $0$ when that set is
empty}; since $\beta_{U_i}(z_i,r_i)>0$ by \eqref{eq:bs-triple}, that set is
nonempty, and it is contained in $(0,1)$ and so bounded above by $1$. The
approximation property is therefore available and gives a parameter
\begin{equation}\label{eq:bs-lam}
  \lambda_i\in(0,1)\ \text{not good for }U_i\text{ at }(z_i,r_i),\qquad
  \lambda_i>\tfrac12\beta_{U_i}(z_i,r_i).
\end{equation}
Combining \eqref{eq:bs-triple} and \eqref{eq:bs-lam}, and using
$\lambda_i\le\beta_{U_i}(z_i,r_i)\le b_d(s_i)$ --- the first because
$\lambda_i$ belongs to the set whose supremum is $\beta_{U_i}(z_i,r_i)$, the
second because $(U_i,z_i,r_i)$ is admissible for $s_i$ --- we obtain
\begin{equation}\label{eq:bs-sandwich}
  \tfrac14b_d(s_i)<\lambda_i\le b_d(s_i).
\end{equation}

\emph{Step 3: $\lambda_i\to0$.}
By \eqref{eq:bs-si}, $s_i\to0$ with $s_i>0$. Let $\varepsilon>0$ and take $s_0$
as in \noderef{FamilyFlatness} for this $\varepsilon$; for all large $i$ we have
$0<s_i\le s_0$ and therefore $b_d(s_i)\le\varepsilon$. Hence $b_d(s_i)\to0$, and
$\lambda_i\le b_d(s_i)$ from \eqref{eq:bs-sandwich} gives
\begin{equation}\label{eq:bs-lamzero}
  \lambda_i\longrightarrow0 .
\end{equation}

\emph{Step 4: the uniformity with $\Lambda=32$.}
Fix $M_*\ge2$. Since $s_i\to0$, $\gamma_i\to0$ and $\lambda_i\to0$, and since
$(2M_*)^{\gamma_i}=\exp(\gamma_i\log(2M_*))\to1$, there is $i_0$ such that for
every $i\ge i_0$
\begin{equation}\label{eq:bs-thresholds}
  M_*s_i<\tfrac12,\qquad 2M_*s_i\le\tfrac18,\qquad
  (2M_*)^{\gamma_i}<4,\qquad 32\lambda_i<1 .
\end{equation}
Let $i\ge i_0$, let $w$ satisfy $|w-z_i|\le M_*r_i$ and let
$0<\rho\le M_*r_i$.

\emph{(4a) $w\in G$.} By \eqref{eq:bs-triple} and \eqref{eq:bs-thresholds},
\begin{equation}\label{eq:bs-Mr}
  M_*r_i\le M_*s_iD(z_i)<\tfrac12D(z_i),
\end{equation}
so $|w-z_i|<\tfrac12D(z_i)$. By \eqref{eq:bs-D}, $|w|\le|z_i|+|w-z_i|
<|z_i|+(\tfrac34-|z_i|)=\tfrac34$, i.e.\ $w\in G$.

\emph{(4b) $\rho\le\tfrac18D(w)$.} Since $D$ is $1$-Lipschitz,
$D(w)\ge D(z_i)-|w-z_i|\ge D(z_i)-M_*r_i\ge\tfrac12D(z_i)$ by
\eqref{eq:bs-Mr}, so $D(z_i)\le2D(w)$ and, using \eqref{eq:bs-triple} and
\eqref{eq:bs-thresholds},
\begin{equation}\label{eq:bs-rho}
  \rho\le M_*r_i\le M_*s_iD(z_i)\le2M_*s_iD(w)\le\tfrac18D(w).
\end{equation}
Thus the center $w$ and radius $\rho$ are admissible.

\emph{(4c) $32\lambda_i$ is good for $U_i$ at $(w,\rho)$.} The triple
$(U_i,w,\rho)$ is admissible for the scale $2M_*s_i$, since $U_i\in\mathcal U$,
$w\in G$ by (4a), $\rho>0$, and $\rho\le2M_*s_iD(w)$ by \eqref{eq:bs-rho}.
Therefore
\[
  \beta_{U_i}(w,\rho)\le b_d(2M_*s_i).
\]
Now $2M_*s_i\ge s_i$ because $2M_*\ge4\ge1$, and $2M_*s_i\le1/8$ by
\eqref{eq:bs-thresholds}; so \eqref{eq:bs-running} applies at $s=2M_*s_i$ and
gives
\[
  b_d(2M_*s_i)\le2b_d(s_i)\frac{(2M_*s_i)^{\gamma_i}}{s_i^{\gamma_i}}
   =2b_d(s_i)(2M_*)^{\gamma_i}.
\]
By \eqref{eq:bs-sandwich}, $b_d(s_i)<4\lambda_i$, and by
\eqref{eq:bs-thresholds}, $(2M_*)^{\gamma_i}<4$; hence
\[
  \beta_{U_i}(w,\rho)\le2b_d(s_i)(2M_*)^{\gamma_i}
   <2\cdot4\lambda_i\cdot4=32\lambda_i .
\]
Finally $32\lambda_i<1$ by \eqref{eq:bs-thresholds}, so
$32\lambda_i\in(\beta_{U_i}(w,\rho),1)$ and therefore $32\lambda_i$ is good for
$U_i$ at $(w,\rho)$: there is $P$ realizing the cylinder fit of parameter
$32\lambda_i$ for $U_i$ at $(w,\rho)$ in the sense of \noderef{IsCylinderFit}.

\emph{Conclusion.}
We check every clause of \noderef{IsBadScaleData} for the data
$(\gamma_i,s_i,r_i,\lambda_i,U_i,z_i)$ just constructed, with $\Lambda=32$.
We have $\gamma_i=1/(i+1)>0$ with $\gamma_i\to0$; $s_i>0$ with $s_i\to0$ by
\eqref{eq:bs-si}; $\lambda_i>0$ by \eqref{eq:bs-lam} with $\lambda_i\to0$ by
\eqref{eq:bs-lamzero}; $U_i\in\mathcal U$ and $z_i\in G$ by Step 2; $r_i>0$ and
$r_i\le s_iD(z_i)$ by \eqref{eq:bs-triple}; $\lambda_i$ is not good for $U_i$ at
$(z_i,r_i)$, i.e.\ no $P$ realizes that cylinder fit, by \eqref{eq:bs-lam}; and
the final clause is Step 4, whose threshold $i_0$ depends on $M_*$ only.
\end{proof}
